\clearpage
\SourceLeaf{83}
\Needspace{10\baselineskip}\begin{center}
\textbf{NEUER BEWEIS DES LAGRANGESCHEN LEHRSATZES.}
\end{center}

\noindent\hspace*{\parindent}Der wichtige Lehrsatz des Herrn \textsc{de la Grange}, nach dem jede Funktion von
$x$ durch eine Reihe nach den Potenzen von $u$ dargestellt wird, wenn $x=t+uX$,
ist schon von seinem Entdecker auf scharfsinnige Weise bewiesen worden\HistoricalNote{*)}{*)}{Hist.\ de l'Ac.\ de Berlin T.\ XXIV, Année 1768 (Berlin 1770), p.\ 275. Die ganze Abhandlung findet man übersetzt im dritten Band, mit dem Hr.\ \textsc{Michelsen Eulers} Analysis des Unendlichen ergänzt hat. Auch hat jüngst Hr.\ D.\ \textsc{Murhard} versucht, diesen Beweis durch einen knappen Auszug in einem eigenen Programm bekannter zu machen.}.
Da dieser Beweis jedoch darauf beruht, dass $X$ die Gestalt
\begin{center}
$\displaystyle ax^2+bx^3+cx^4+.\,.\,.\,.$\HistoricalNote{**)}{**)}{Dasselbe gilt für die Umkehrungsformel, die Hr.\ \textsc{Eschenbach} unabhängig gefunden und Hr.\ M.\ \textsc{Rothe} bewiesen hat und die sich leicht mit dem Lehrsatz von \textsc{Lagrange} verbinden lässt.}
\end{center}
\noindent hat, und zudem recht weitschweifig ist, haben verschiedene Geometer versucht,
den Satz allein aus der Natur der Funktionen herzuleiten\HistoricalNote{***)}{***)}{\textsc{Lexell}, Comm.\ nov.\ Petr., T.\ XVI (1772), p.\ 230.\par
\textsc{Condorcet}, Miscell.\ Taurin., T.\ V (1770--73), Classe Math.\ p.\ 7.\par
\textsc{Cousin}, Astronomie physique, p.\ 15 (Paris 1787).\par
\textsc{Pfaff}, Archiv der r.\ u.\ a.\ M., 1.\ Band (1795), S.\ 81.}.
Von all diesen Bemühungen ist jedoch nur der Beweis von Herrn\ Prof.\ \textsc{Pfaff}
völlig streng. Die übrigen sind im Grunde bloße Induktion. Ich hoffe,
mancher Leser wird hier gerne einen weiteren, ebenfalls strengen Beweis
finden, der mit dem von Herrn\ Prof.\ \textsc{Pfaff} nichts gemein hat. Die
Hauptidee ist im Wesentlichen dieselbe wie bei \textsc{Cousin}, doch die
Ausführung scheint mir neu und einfach. Unterschiedliche Methoden tragen
immer dazu bei, eine Wahrheit besser zu beleuchten, und der vorliegende
Satz scheint bei uns bislang weniger bekannt zu sein, als er es verdient.

